% Generated draft: requires mathematical and source review.
\documentclass[11pt]{article}
\usepackage[margin=25mm]{geometry}
\usepackage{fontspec}
\setmainfont{Latin Modern Roman}
\usepackage{amsmath,amssymb,mathtools,unicode-math}
\setmathfont{Latin Modern Math}
\usepackage{enumitem,needspace,xurl,hyperref}
\setlength{\parindent}{0pt}
\setlength{\parskip}{5pt}
\setlength{\emergencystretch}{3em}
\newcommand{\N}{\mathbb{N}}
\newcommand{\Z}{\mathbb{Z}}
\newcommand{\Q}{\mathbb{Q}}
\newcommand{\R}{\mathbb{R}}
\newcommand{\C}{\mathbb{C}}
\newcommand{\F}{\mathbb{F}}
\newcommand{\A}{\mathbb{A}}
\newcommand{\PP}{\mathbb{P}}
\newcommand{\E}{\mathbb{E}}
\newcommand{\eps}{\varepsilon}
\newcommand{\dd}{\,\mathrm{d}}
\begin{document}
\section*{UnsolvedMath ID: 1200022}
\textbf{Source problem code:} catalogue entry
\subsection{Definitions and mathematical statement}
A unimodular random network is a probability law $\mu$ on isomorphism classes of rooted connected locally finite networks $(T,o)$ satisfying the mass-transport principle: for every nonnegative measurable isomorphism-invariant function $f(T,u,v)$,
\[
\mathbb E_\mu\sum_{v\in V(T)}f(T,o,v)
=\mathbb E_\mu\sum_{v\in V(T)}f(T,v,o).
\]
It is ergodic if each measurable event invariant under changing the root within its connected network has probability $0$ or $1$. The hypothesis in Ab\'ert's question is that $T$ is almost surely an infinite tree.

Under the standard bounded-degree unmarked interpretation, an expander family consists of finite graphs $(G_j)$ with $|V(G_j)|\to\infty$, a common maximum-degree bound $D<\infty$, and a constant $h>0$ such that
\[
|\partial_E S|\geq h|S|
\quad\text{whenever }\varnothing\ne S\subseteq V(G_j),\quad |S|\leq\tfrac12|V(G_j)|,
\]
for every $j$, where $\partial_E S$ is the set of edges joining $S$ to its complement. If $U_j$ is a uniform random vertex of $G_j$, local weak convergence to $\mu$ means that for every radius $r$ the distributions of the rooted radius-$r$ neighborhoods of $(G_j,U_j)$ converge to that of $(T,o)$.

The source asks whether every ergodic unimodular law $\mu$ almost surely supported on infinite trees can be represented in this way by an expander family. The definition of an expander family and the meaning of ``network'' and ``limit'' are not specified in the supplied text. The displayed definitions give a common interpretation, not a verified restriction of the source to bounded-degree unmarked networks. In particular, the literal target class need not have bounded degrees or any stated nonamenability condition and includes amenable trees such as the bi-infinite path. The record does not explain the compatibility of that scope with a uniform positive expansion requirement; these hypotheses must be checked before accepting the universal assertion as the intended open target.

The supplied summary labels the converse realization question OPEN-TRIAGE, unresolved as far as its author could verify. It also says the forward implication from expanders to ``such limits'' is established, but supplies no theorem or hypotheses. General expansion alone does not impose a tree local limit; a tree-like or girth hypothesis for that part is missing from the supplied summary.
\subsection{Short English statement}
The question asks whether every ergodic unimodular random infinite tree can arise as a local limit of finite expanders. The source does not specify degree bounds, marks, or the expansion convention, and its literal class includes amenable trees. Those scope issues need resolution before the reported open classification can be interpreted precisely.
\subsection{Sources}
\begin{itemize}
\item Local UnsolvedMath registry, permanent ID 1200022, exact source code catalogue entry, ``Some Questions --- Question 22''; status open, classification OPEN-TRIAGE.
\item Ab\'ert, \emph{Some questions} (2010), Question 22, \url{https://www.renyi.hu/~abert/questions.pdf}. The provenance records manual transcription of a public author-hosted PDF, accessed July 29, 2026. It says the 2010 list was inspected, while current status was not independently verified.
\item The supplied research summary reports the realization question unresolved as far as its author could verify and calls the forward direction established. No independent authoritative citation, theorem statement, or supporting hypotheses for either assessment are supplied.
\end{itemize}
\end{document}
